% !TeX root = part01.tex

\documentclass[../final.tex]{subfiles}

\begin{document}

% ============================================================
% Лекция 01.09
% Представление чисел в ЭВМ
% ============================================================

\chapter{Представление чисел в ЭВМ}

\rsubtitle{Двоичная арифметика и машинная точность}


% ============================================================
\section{Целые числа}
% ============================================================

\subsection{Двоичное кодирование}

В ЭВМ вся информация хранится в виде последовательности
\rtxt{3}{битов}. Каждый бит может принимать только два значения:
%
\[
    0
    \qquad \text{или} \qquad
    1.
\]

Следовательно, с помощью $p$ бит можно закодировать
%
\[
    2^p
\]
%
различных последовательностей.

Если интерпретировать последовательность
%
\[
    a_{p-1}a_{p-2}\ldots a_1a_0,
    \qquad
    a_i\in\{0,1\},
\]
%
как двоичное число, то её значение равно
%
\begin{equation*}
    A
    =
    a_{p-1}2^{p-1}
    +
    a_{p-2}2^{p-2}
    +\ldots+
    a_1 2
    +
    a_0
    =
    \sum_{i=0}^{p-1} a_i2^i.
\end{equation*}

Таким способом кодируются беззнаковые целые числа
%
\[
    0,1,\ldots,2^p-1.
\]

Например, для $p=32$ имеется
%
\[
    2^{32}
\]
%
различных битовых последовательностей.

Для стандартного 32-битного целого числа
%
\[
    4\ \text{байта}\times8\ \text{бит}
    =
    32\ \text{бита}.
\]


% ------------------------------------------------------------
\subsection{Отрицательные числа}
% ------------------------------------------------------------

Для представления целых чисел со знаком необходимо закодировать
как положительные, так и отрицательные значения.

В дополнительном коде отрицательное число $-N$ представляется
числом
%
\begin{equation*}
    2^p-N.
\end{equation*}

Иными словами, код числа $-N$ выбирается так, чтобы
%
\begin{equation*}
    N + (2^p-N)
    =
    2^p
    \equiv
    0
    \pmod{2^p}.
\end{equation*}

Поэтому арифметику $p$-битных целых чисел удобно рассматривать
как арифметику по модулю $2^p$.

Например,
%
\begin{equation*}
    -1
    \longleftrightarrow
    2^p-1
    =
    \underbrace{11\ldots11}_{p\text{ бит}}{}_2 .
\end{equation*}

Диапазон знаковых $p$-битных целых чисел имеет вид
%
\begin{equation*}
    -2^{p-1}
    \leq
    N
    \leq
    2^{p-1}-1.
\end{equation*}

\textit{Замечание.}
Именно поэтому при переполнении целочисленной переменной
результат фактически ``переходит через границу'' доступного
диапазона.


% ============================================================
\section{Числа с плавающей точкой}
% ============================================================

\subsection{Фиксированная и плавающая точка}

Вещественное число можно хранить в форме с
\rtxt{3}{фиксированной точкой}, заранее задавая положение
десятичного разделителя.

Например,
%
\[
    1{,}031
    =
    \frac{1031}{1000},
    \qquad
    10{,}05
    =
    \frac{1005}{100}.
\]

Такой способ неудобен для чисел, изменяющихся в очень широком
диапазоне. Поэтому на практике используется
\rtxt{3}{форма с плавающей точкой}:
%
\begin{equation*}
    x
    =
    M\cdot b^E,
\end{equation*}
%
где $M$ --- мантисса, $E$ --- порядок, а $b$ --- основание
системы счисления.

Например, в десятичной системе
%
\[
    1{,}031
    =
    1{,}031\cdot10^0,
\]
%
\[
    10{,}51
    =
    1{,}051\cdot10^1.
\]


% ------------------------------------------------------------
\subsection{Нормализованная форма}
% ------------------------------------------------------------

Чтобы запись числа с плавающей точкой была однозначной,
используется \rtxt{3}{нормализованная форма}: положение запятой
выбирается так, чтобы перед ней находилась ровно одна ненулевая
цифра.

Например,
%
\[
    1{,}031\cdot10^0
\]
%
является нормализованной записью, тогда как
%
\[
    10{,}31\cdot10^{-1}
\]
%
описывает то же число, но нормализованной записью не является.

В двоичной системе первая цифра нормализованной мантиссы всегда
равна единице. Поэтому число можно записать в форме
%
\begin{equation*}
    x
    =
    (-1)^S
    \left(
        1+\frac{M}{2^n}
    \right)
    2^E,
\end{equation*}
%
где
%
\[
    S\in\{0,1\},
    \qquad
    M=0,1,\ldots,2^n-1.
\]

Здесь

\begin{itemize}
    \item $S$ --- бит знака;
    \item $M$ --- целое число, кодирующее дробную часть мантиссы;
    \item $n$ --- число бит дробной части мантиссы;
    \item $E$ --- порядок.
\end{itemize}

Поскольку старшая единица нормализованной мантиссы всегда известна,
её нет необходимости явно хранить в памяти. Это так называемая
\rtxt{3}{скрытая единица}.


% ------------------------------------------------------------
\subsection{Формат IEEE 754}
% ------------------------------------------------------------

В памяти нормализованное число записывается в виде трёх полей:
%
\[
    \boxed{S}
    \qquad
    \boxed{E_S}
    \qquad
    \boxed{M},
\]
%
где $S$ --- знак, $E_S$ --- смещённый порядок, $M$ --- дробная
часть мантиссы.

Если под порядок выделено $w$ бит, вводится смещение
%
\begin{equation*}
    B
    =
    2^{w-1}-1.
\end{equation*}

Тогда вместо самого порядка $E$ хранится
%
\begin{equation*}
    E_S
    =
    E+B
    =
    E+2^{w-1}-1.
\end{equation*}

Следовательно,
%
\begin{equation*}
    E
    =
    E_S-\left(2^{w-1}-1\right),
\end{equation*}
%
и нормализованное число принимает вид
%
\begin{equation*}
    x
    =
    (-1)^S
    \left(
        1+\frac{M}{2^n}
    \right)
    2^{\,E_S-(2^{w-1}-1)}.
\end{equation*}

Два значения поля порядка зарезервированы для специальных случаев:
%
\[
    E_S=0,
    \qquad
    E_S=2^w-1.
\]

Поэтому для нормализованных чисел
%
\begin{equation*}
    1
    \leq
    E_S
    \leq
    2^w-2.
\end{equation*}

Подставляя
%
\[
    E_S=E+2^{w-1}-1,
\]
%
получаем
%
\begin{align*}
    1
    &\leq
    E+2^{w-1}-1
    \leq
    2^w-2,
    \\[3pt]
    -2^{w-1}+2
    &\leq
    E
    \leq
    2^{w-1}-1.
\end{align*}

Таким образом,
%
\begin{equation*}
    E_{\min}
    =
    -2^{w-1}+2,
    \qquad
    E_{\max}
    =
    2^{w-1}-1.
\end{equation*}


% ------------------------------------------------------------
\subsection{Одинарная точность: \texttt{float}}
% ------------------------------------------------------------

Для формата IEEE 754 binary32, которому обычно соответствует
тип \texttt{float},
%
\[
    32
    =
    1+8+23.
\]

Биты распределяются следующим образом:
%
\[
    \underbrace{\boxed{S}}_{1\text{ бит}}
    \qquad
    \underbrace{\boxed{E_S}}_{8\text{ бит}}
    \qquad
    \underbrace{\boxed{M}}_{23\text{ бита}}.
\]

То есть
%
\[
    w=8,
    \qquad
    n=23,
    \qquad
    B=2^7-1=127.
\]

Отсюда
%
\begin{equation*}
    E_{\min}
    =
    -126,
    \qquad
    E_{\max}
    =
    127.
\end{equation*}

Минимальное положительное нормализованное число получается при
$M=0$ и $E=E_{\min}$:
%
\begin{equation*}
    f_0
    =
    2^{E_{\min}}
    =
    2^{-126}
    \approx
    1.18\cdot10^{-38}.
\end{equation*}


Следующее машинное число имеет $M=1$:
%
\begin{equation*}
    f_1
    =
    \left(
        1+\frac{1}{2^n}
    \right)
    2^{E_{\min}}.
\end{equation*}

Расстояние между ними равно
%
\begin{equation*}
    f_1-f_0
    =
    2^{E_{\min}-n}.
\end{equation*}

При этом
%
\begin{equation*}
    \frac{f_0}{f_1-f_0}
    =
    2^n.
\end{equation*}

Для \texttt{float}, где $n=23$,
%
\[
    2^{23}
    \approx
    8.4\cdot10^6.
\]

То есть без дополнительного механизма около нуля возникла бы
очень большая область непредставимых чисел.


% ------------------------------------------------------------
\subsection{Денормализованные числа}
% ------------------------------------------------------------

Чтобы заполнить область около нуля, вводятся
\rtxt{3}{денормализованные} (subnormal) числа.

Для них поле порядка полностью заполнено нулями:
%
\[
    E_S=0.
\]

Скрытая единица в мантиссе при этом отсутствует, поэтому число
записывается как
%
\begin{equation*}
    x
    =
    (-1)^S
    \frac{M}{2^n}
    2^{E_{\min}}.
\end{equation*}

Минимальное положительное денормализованное число соответствует
$M=1$:
%
\begin{equation*}
    f_{\min}^{\mathrm{denorm}}
    =
    2^{E_{\min}-n}.
\end{equation*}

Максимальное денормализованное число:
%
\begin{equation*}
    f_{\max}^{\mathrm{denorm}}
    =
    \left(
        1-\frac{1}{2^n}
    \right)
    2^{E_{\min}}.
\end{equation*}

А следующее за ним число уже является минимальным
нормализованным:
%
\begin{equation*}
    f_{\min}^{\mathrm{norm}}
    =
    2^{E_{\min}}.
\end{equation*}

Таким образом, переход от денормализованных чисел к
нормализованным происходит без большого разрыва.


% ------------------------------------------------------------
\subsection{Специальные значения}
% ------------------------------------------------------------

Зарезервированные значения поля порядка позволяют кодировать
ноль, бесконечность и \texttt{NaN}.

Для краткости обозначим через $00\ldots0$ последовательность из нулей,
а через $11\ldots1$ --- последовательность из единиц.

\begin{enumerate}
    \item
    \[
        \boxed{S}\;
        \boxed{00\ldots0}\;
        \boxed{00\ldots0}
    \]
    соответствует числу $+0$ или $-0$ в зависимости от $S$.

    \item
    \[
        \boxed{S}\;
        \boxed{00\ldots0}\;
        \boxed{M},
        \qquad
        M\neq0,
    \]
    соответствует денормализованному числу.

    \item
    \[
        \boxed{S}\;
        \boxed{11\ldots1}\;
        \boxed{00\ldots0}
    \]
    соответствует
    \[
        +\infty
        \qquad\text{или}\qquad
        -\infty.
    \]

    \item
    \[
        \boxed{S}\;
        \boxed{11\ldots1}\;
        \boxed{M},
        \qquad
        M\neq0,
    \]
    соответствует \rtxt{3}{NaN} (\textit{Not a Number}).

    \item При
    \[
        1\leq E_S\leq2^w-2
    \]
    кодируется обычное нормализованное число.
\end{enumerate}


% ------------------------------------------------------------
\subsection{Машинная точность}
% ------------------------------------------------------------

Поскольку число бит мантиссы конечно, результат арифметической
операции в общем случае приходится округлять до ближайшего
представимого машинного числа.

Например,
%
\[
    1{,}08\cdot2{,}03
    =
    2{,}1924.
\]

Если оставить только три значащие цифры, получим
%
\[
    2{,}1924
    \longrightarrow
    2{,}19,
\]
%
что вносит относительную погрешность порядка
%
\[
    10^{-3}.
\]

Рассмотрим теперь машинные числа около единицы.
При $E=0$ они имеют вид
%
\begin{equation*}
    f
    =
    1+\frac{M}{2^n}.
\end{equation*}

Два соседних числа:
%
\begin{equation*}
    f_0=1,
    \qquad
    f_1
    =
    1+\frac{1}{2^n}.
\end{equation*}

Расстояние между ними равно
%
\begin{equation*}
    f_1-f_0
    =
    \frac{1}{2^n}.
\end{equation*}

Эта величина называется \rtxt{3}{машинным эпсилон}:
%
\begin{equation*}
    \boxed{
        \varepsilon
        =
        2^{-n}
    }.
\end{equation*}

Она характеризует относительное расстояние между соседними
машинными числами порядка единицы.

При округлении к ближайшему числу максимальная относительная
погрешность имеет порядок
%
\begin{equation*}
    \frac{\varepsilon}{2}
    =
    2^{-(n+1)}.
\end{equation*}

Практически машинное эпсилон можно характеризовать условиями
%
\begin{align*}
    1+\varepsilon
    &\neq
    1,
    \\[3pt]
    1+\frac{\varepsilon}{2}
    &=
    1
    \qquad
    \text{в машинной арифметике}.
\end{align*}

Для одинарной точности
%
\begin{equation*}
    \varepsilon_{\texttt{float}}
    =
    2^{-23}
    \approx
    1.19\cdot10^{-7},
\end{equation*}
%
а для двойной точности
%
\begin{equation*}
    \varepsilon_{\texttt{double}}
    =
    2^{-52}
    \approx
    2.22\cdot10^{-16}.
\end{equation*}

\end{document}